\metatitle{Generalizations of permutations}
\metadescription{Affine permutations, signed permutations, colored permutations, and Stirling permutations.}


\section[notationAffinePermutations]{Affine permutations}

The set of \defin{affine permutations} $\symAff_n$ is the set of bijections $\pi: \setZ \to \setZ$
satisfying $\pi(i+n)=\pi(i)+n$ for all $i\in \setZ$ and $\pi(1)+\pi(2)+\dotsb+\pi(n)=1+2+\dotsb+n$.


\section[permutationsCayley]{Cayley permutations}

A \defin{Cayley permutation} of length $n$ is a word
$w=(w_1,\dotsc,w_n)$ of positive integers whose set of values is $[k]$ for
some $k$.  Equivalently, these are packed words.

\name{Giulio Cerbai} and \name{Anders Claesson} count fixed-point-free Cayley
permutations using two-sort species and differential equations for the
corresponding functional digraphs \cite{CerbaiClaesson2025x}.  They obtain an
explicit formula, prove that the fixed-point-free proportion tends to $1/e$,
and also count the cases where the functional digraph is a tree, forest, or
connected.


\section[permutationsTypeB]{Permutations of type $B$}

The set of permutations of type $B$, is defined as the permutations on $B_n \coloneqq \{\pm 1, \pm 2,\dotsc, \pm n\}$
with the property that $\pi(-i)=-\pi(i)$. This is also known as the \hyperref[weylLieGroups]{hyperoctahedral group}.

\name{Eli Bagno}, \name{Riccardo Biagioli}, \name{Frédéric Jouhet}, and
\name{Yuval Roichman} study fully commutative elements in type $B$ through
block number and descents \cite{BagnoBiagioliJouhetRoichman2022}.  They prove
a Schur-positivity result for the corresponding descent enumerators.


\section[permutationsColored]{Colored permutations}

A \defin{$r$-colored permutation} of size $n$ is a permutation in which each
letter is assigned one of $r$ colors.  Equivalently, it is a pair
\[
  (\pi,c) \in \symS_n \times \{0,1,\dotsc,r-1\}^n,
\]
usually written in colored window notation as
\[
  \pi(1)^{(c_1)}\,\pi(2)^{(c_2)}\dotsm \pi(n)^{(c_n)}.
\]
The set of all such objects is the wreath product
\[
  C_r \wr \symS_n \cong (\setZ/r\setZ)^n \rtimes \symS_n.
\]
For example,
\[
  3^{(2)}\,1^{(0)}\,2^{(1)}
\]
is a $3$-colored permutation of size $3$.

When $r=1$ we recover ordinary \hyperref[notationPermutations]{permutations};
when $r=2$ we recover the signed permutations, or type $B$ permutations, after
identifying the two colors with signs.  Many permutation statistics, including
descents and major index, have colored analogues; see
\cite{GarsiaGessel1979,Hyatt2012}.  Colored permutations also appear in
wreath-product Schur theory and in ribbon-tableau versions of
\hyperref[rsk]{RSK}.


\section[permutationsStirling]{Stirling permutations}

In \cite{GesselStanley1978}, \name[Ira Gessel]{Gessel} and \name[Richard P. Stanley]{Stanley}
introduced the set of \defin{Stirling permutations, $Q_n$}.
These are permutations of the multiset $\{1,1,2,2,3,3,\dotsc,n,n\}$ with the additional property
that for between the two entries equal to $i \in [n]$, only entries larger than $i$ may appear.
For example, $12234431$ is an element in $Q_n$.

We have that $|Q_n| = (2n-1)!!$.
